% ============================================================
% paper/sections/00a_epistemic_status_citation_contract.tex
% FWHQ-02 — Epistemic Positioning & Publication Hardening
% Scope: extension-only; no new theory; no empirical claims
% ============================================================

\section{Epistemic Status of This Artifact}
\label{sec:epistemic-status}

\subsection{Artifact classification}
This document is a human-readable projection of an \emph{immutable,
executable-closed formal specification}, namely
\texttt{ETS.K\_EA.v1.1}.

Its epistemic category is that of a \textbf{formal object and diagnostic
calculus}. It defines symbols, admissible configurations, constraints,
boundary conditions (including STOP), and closure properties.
All statements are to be read as \emph{intra-formal}: they hold
\emph{within the closed ETS-defined system} and carry no claims beyond it.

\subsection{What kind of knowledge this is}
The knowledge encoded in this artifact is:
\begin{itemize}
  \item \textbf{Definitional}: it fixes the meaning of symbols and operators.
  \item \textbf{Structural}: it specifies relations, constraints, and
        admissibility conditions.
  \item \textbf{Negative and boundary-aware}: it explicitly encodes
        non-existence, STOP, and undefinedness as first-class outcomes.
\end{itemize}

It is knowledge of \emph{form and constraint}, not of empirical fact,
prediction, or performance.

\subsection{What this artifact is not}
To eliminate category errors, this artifact is explicitly \textbf{not}:
\begin{itemize}
  \item a scientific theory making empirical or testable claims;
  \item an empirical study, benchmark, or validation report;
  \item a methodology, framework, or best-practice prescription;
  \item a comparative positioning against other EA approaches;
  \item a claim of usefulness, effectiveness, resilience, or success.
\end{itemize}

Requests for empirical grounding, case studies, comparisons, validation,
or usefulness justification are therefore \emph{out-of-scope by document
type}, not missing content.

\subsection{Why this artifact exists as a publication}
This artifact exists as a citable publication in order to provide a
\textbf{stable, versioned, and inspectable reference} for a closed formal
object.

Its publication purpose is:
\begin{itemize}
  \item to make the formal specification addressable by reference;
  \item to allow unambiguous citation of definitions and constraints;
  \item to provide an immutable anchor for executable closure and STOP semantics.
\end{itemize}

Publication does not imply empirical authority, validation, or endorsement.
It serves archival, referential, and infrastructural functions only.

\section{Citation \& Use Contract}
\label{sec:citation-contract}

\subsection{What a citation to this artifact asserts}
A valid citation to this document asserts one or more of the following:
\begin{enumerate}
  \item the existence and identity of a specific, versioned formal object;
  \item the definitions, symbols, constraints, and boundary semantics
        introduced therein;
  \item the executable-closed nature of the specification at the level of
        file identity and reconstructability.
\end{enumerate}

\subsection{What a citation does not assert}
A citation to this artifact does \textbf{not} assert:
\begin{itemize}
  \item empirical correctness or validation;
  \item predictive accuracy or explanatory power about real systems;
  \item practical usefulness or effectiveness;
  \item superiority over alternative approaches;
  \item any real-world outcome or guarantee.
\end{itemize}

\subsection{Invalid citation and use patterns}
The following uses constitute category errors and are invalid by definition:
\begin{itemize}
  \item citing this artifact as empirical evidence;
  \item citing it as prescriptive guidance or recommendation;
  \item citing it as a comparative benchmark;
  \item treating publication or archiving as proof of scientific validity.
\end{itemize}

\subsection{Versioning and pinning requirements}
Formal use of this artifact requires citation of a \textbf{specific version}
(release, tag, or commit) of the document.
Where arguments depend on closure properties or STOP semantics, the citation
must be pinned to the corresponding ETS identity anchor stated in the front
matter.
